\documentclass[10pt,a4paper]{amsart}
\usepackage{amsmath,amssymb,amsthm,graphicx}
% For commutative diagrams we use Xy-pic
\usepackage[all]{xy}

% We use 2cell for 2-commutative diagrams.
\xyoption{2cell}
% Theorem environments.
%
\theoremstyle{plain}
\newtheorem{theorem}[subsection]{Theorem}
\newtheorem{proposition}[subsection]{Proposition}
\newtheorem{lemma}[subsection]{Lemma}

\theoremstyle{definition}
\newtheorem{definition}[subsection]{Definition}
\newtheorem{example}[subsection]{Example}
\newtheorem{exercise}[subsection]{Exercise}
\newtheorem{situation}[subsection]{Situation}

\theoremstyle{remark}
\newtheorem{remark}[subsection]{Remark}
\newtheorem{remarks}[subsection]{Remarks}

\numberwithin{equation}{subsection}

% Macros
%
\def\lim{\mathop{\mathrm{lim}}\nolimits}
\def\colim{\mathop{\mathrm{colim}}\nolimits}
\def\Spec{\mathop{\mathrm{Spec}}}
\def\Hom{\mathop{\mathrm{Hom}}\nolimits}
\def\Ext{\mathop{\mathrm{Ext}}\nolimits}
\def\SheafHom{\mathop{\mathcal{H}\!\mathit{om}}\nolimits}
\def\SheafExt{\mathop{\mathcal{E}\!\mathit{xt}}\nolimits}
\def\Sch{\mathit{Sch}}
\def\Mor{\mathop{\mathrm{Mor}}\nolimits}
\def\Ob{\mathop{\mathrm{Ob}}\nolimits}
\def\Sh{\mathop{\mathit{Sh}}\nolimits}
\def\NL{\mathop{N\!L}\nolimits}
\def\CH{\mathop{\mathrm{CH}}\nolimits}
\def\proetale{{pro\text{-}\acute{e}tale}}
\def\etale{{\acute{e}tale}}
\def\QCoh{\mathit{QCoh}}
\def\Ker{\mathop{\mathrm{Ker}}}
\def\Im{\mathop{\mathrm{Im}}}
\def\Coker{\mathop{\mathrm{Coker}}}
\def\Coim{\mathop{\mathrm{Coim}}}

% Boxtimes
%
\DeclareMathSymbol{\boxtimes}{\mathbin}{AMSa}{"02}

%
% Macros for moduli stacks/spaces
%
\def\QCohstack{\mathcal{QC}\!\mathit{oh}}
\def\Cohstack{\mathcal{C}\!\mathit{oh}}
\def\Spacesstack{\mathcal{S}\!\mathit{paces}}
\def\Quotfunctor{\mathrm{Quot}}
\def\Hilbfunctor{\mathrm{Hilb}}
\def\Curvesstack{\mathcal{C}\!\mathit{urves}}
\def\Polarizedstack{\mathcal{P}\!\mathit{olarized}}
\def\Complexesstack{\mathcal{C}\!\mathit{omplexes}}
% \Pic is the operator that assigns to X its picard group, usage \Pic(X)
% \Picardstack_{X/B} denotes the Picard stack of X over B
% \Picardfunctor_{X/B} denotes the Picard functor of X over B
\def\Pic{\mathop{\mathrm{Pic}}\nolimits}
\def\Picardstack{\mathcal{P}\!\mathit{ic}}
\def\Picardfunctor{\mathrm{Pic}}
\def\Deformationcategory{\mathcal{D}\!\mathit{ef}}
\newenvironment{reference}{}{}
\newenvironment{slogan}{}{}
\newenvironment{history}{}{}
\usepackage{cleveref}
\crefname{originallemma}{original lemma}{original lemmas}
\Crefname{originallemma}{Original lemma}{Original lemmas}
\makeatletter
\newlabel{lemma-divisor-special-fiber}{{0C5Z}{source}{}{sourcepointer-1}{}}
\newlabel{lemma-divisor-special-fiber@cref}{{[originallemma][0][]0C5Z}{[1][0][]source}{}{}{}}
\makeatother
\crefname{originallemma}{original lemma}{original lemmas}
\Crefname{originallemma}{Original lemma}{Original lemmas}
\makeatletter
\newlabel{lemma-regular-model-connected}{{0C62}{source}{}{sourcepointer-2}{}}
\newlabel{lemma-regular-model-connected@cref}{{[originallemma][0][]0C62}{[1][0][]source}{}{}{}}
\makeatother
\crefname{originallemma}{original lemma}{original lemmas}
\Crefname{originallemma}{Original lemma}{Original lemmas}
\makeatletter
\newlabel{lemma-properties-form}{{0C66}{source}{}{sourcepointer-3}{}}
\newlabel{lemma-properties-form@cref}{{[originallemma][0][]0C66}{[1][0][]source}{}{}{}}
\makeatother
\crefname{originallemma}{original lemma}{original lemmas}
\Crefname{originallemma}{Original lemma}{Original lemmas}
\makeatletter
\newlabel{divisors-lemma-ses-add-divisor}{{0C4T}{source}{}{sourcepointer-4}{}}
\newlabel{divisors-lemma-ses-add-divisor@cref}{{[originallemma][0][]0C4T}{[1][0][]source}{}{}{}}
\makeatother
\crefname{originallemma}{original lemma}{original lemmas}
\Crefname{originallemma}{Original lemma}{Original lemmas}
\makeatletter
\newlabel{varieties-lemma-proper-geometrically-reduced-global-sections}{{0BUG}{source}{}{sourcepointer-5}{}}
\newlabel{varieties-lemma-proper-geometrically-reduced-global-sections@cref}{{[originallemma][0][]0BUG}{[1][0][]source}{}{}{}}
\makeatother
\crefname{originallemma}{original lemma}{original lemmas}
\Crefname{originallemma}{Original lemma}{Original lemmas}
\makeatletter
\newlabel{varieties-lemma-check-invertible-sheaf-trivial}{{0B40}{source}{}{sourcepointer-6}{}}
\newlabel{varieties-lemma-check-invertible-sheaf-trivial@cref}{{[originallemma][0][]0B40}{[1][0][]source}{}{}{}}
\makeatother
\crefname{originallemma}{original lemma}{original lemmas}
\Crefname{originallemma}{Original lemma}{Original lemmas}
\makeatletter
\newlabel{algebra-lemma-integral-under-field}{{00GR}{source}{}{sourcepointer-7}{}}
\newlabel{algebra-lemma-integral-under-field@cref}{{[originallemma][0][]00GR}{[1][0][]source}{}{}{}}
\makeatother
\title{The chosen divisor leads to a chain of subschemes of the special fibre whose global sections are the residue field}
\author{The Stacks Project Authors}
\date{stacks-project:models.tex:local-snapshot-20261012; GNU FDL 1.2 or later}
\begin{document}
\maketitle
\noindent\textit{Editorial note.} The counted unit is the Stacks Project Lemma lemma-regular-model-field at source lines 4812-4826 together with its proof at source lines 4828-4865, a straight adjacent pairing in the chapter file. The chapter is the upstream file models.tex; the pool directory carries the longer name stacks-models and is not used as the chapter name anywhere in this card.\par\medskip
\noindent\textit{Editorial note.} The statement opens with the author's reference environment, which in this chapter is a comment environment; the retained wrapper makes it a visible no-op for this lane, and the environment quotes the cited article on degenerate fibres and stable reduction with an optional argument naming a numbered lemma of that article. That citation is the only one in this card, and its entry is supplied as a task-directory fragment holding it exactly as a bibtex run on the host had formatted it out of the Stacks Project's shared bibliography; the bibliography backend is therefore none.\par\medskip
\noindent\textit{Editorial note.} One context span carries the regular-model situation at lines 4628-4637, which the retained statement names and which its own reference to the divisor lemma makes resolvable through a pointer; it is context rather than a counted unit and does not overlap the counted interval.\par\medskip
\noindent\textit{Editorial note.} The retained proof names three lemmas of this chapter and three results of the Divisors, Varieties and Algebra chapters. All six are declared as non-load-bearing navigation pointers with their stable Stacks tags looked up in the shared tag index; the first three are counted by cards of this lane and none of their intervals is reproduced here.\par\medskip
\noindent\textit{Editorial note.} The retained proof span starts at line 4828, the line of the opening marker read from the source, and contains exactly one proof environment.\par\medskip
\noindent\textit{Allowed standard prerequisites (editorial).}\begin{itemize}
\item The regular-model situation, the divided special fibre and the residue field of a component. The situation fixes a discrete valuation ring, a smooth projective curve over its fraction field, a regular proper model and the components of the special fibre with their multiplicities; dividing the special fibre by the greatest common divisor of the multiplicities yields an effective Cartier divisor, and its multiples are again effective Cartier divisors on the model. The special fibre is connected, each of its components is proper and geometrically reduced over the residue field, and the global sections of a proper geometrically reduced scheme over a field form a field extension; a line bundle on a proper curve whose restriction to each component is trivial is trivial. The author's own lemmas of the chapter and the quoted results of the Algebra, Divisors and Varieties chapters are the allowed inputs.. The counted lemma is the technical heart of the section: it produces the chain of thickenings of the special fibre that the following lemmas use to compute cohomology, and the earlier card on the first cohomology of the reduction points at it.
\end{itemize}\medskip
\ifdefined\hypertarget\hypertarget{sourcepointer-1}{}\fi
\noindent\textit{Source reference.} Original lemma 0C5Z: focused-seeds/AGP/sources/stacks-snapshot-20261011/tags.txt, source lines 15279--15279. Navigation pointer carried by the retained situation, which names this lemma when it writes the special fibre as a sum of its components with multiplicities, and invoked again by the retained proof; that lemma is counted by a separate card of this lane, so its proof is not reproduced here. The number is the stable Stacks tag of the target, looked up in the shared tag index.\par\medskip
\ifdefined\hypertarget\hypertarget{sourcepointer-2}{}\fi
\noindent\textit{Source reference.} Original lemma 0C62: focused-seeds/AGP/sources/stacks-snapshot-20261011/tags.txt, source lines 15282--15282. Navigation pointer invoked by the retained proof for the connectedness of the special fibre, which the chain construction uses; that lemma is counted by a separate card of this lane, so its proof is not reproduced here. The number is the stable Stacks tag of the target, looked up in the shared tag index.\par\medskip
\ifdefined\hypertarget\hypertarget{sourcepointer-3}{}\fi
\noindent\textit{Source reference.} Original lemma 0C66: focused-seeds/AGP/sources/stacks-snapshot-20261011/tags.txt, source lines 15286--15286. Navigation pointer invoked by the retained proof for the numerical properties of the intersection form of a numerical type; the cited statement is an allowed input of the same chapter and its proof is not reproduced here. The number is the stable Stacks tag of the target, looked up in the shared tag index.\par\medskip
\ifdefined\hypertarget\hypertarget{sourcepointer-4}{}\fi
\noindent\textit{Source reference.} Original lemma 0C4T: focused-seeds/AGP/sources/stacks-snapshot-20261011/tags.txt, source lines 15238--15238. Allowed input of the Divisors chapter quoted by the retained proof for the short exact sequence comparing a divisor with the sum of its summands, which is what produces the successive quotients of the chain; the pointer is registered so the reference resolves and the cited statement is not reproduced inside the span. The number is the stable Stacks tag of the target, looked up in the shared tag index.\par\medskip
\ifdefined\hypertarget\hypertarget{sourcepointer-5}{}\fi
\noindent\textit{Source reference.} Original lemma 0BUG: focused-seeds/AGP/sources/stacks-snapshot-20261011/tags.txt, source lines 14874--14874. Allowed input of the Varieties chapter quoted by the retained proof for the fact that the global sections of a proper geometrically reduced scheme over a field form a field extension; the pointer is registered so the reference resolves and the cited statement is not reproduced inside the span. The number is the stable Stacks tag of the target, looked up in the shared tag index.\par\medskip
\ifdefined\hypertarget\hypertarget{sourcepointer-6}{}\fi
\noindent\textit{Source reference.} Original lemma 0B40: focused-seeds/AGP/sources/stacks-snapshot-20261011/tags.txt, source lines 13983--13983. Allowed input of the Varieties chapter quoted by the retained proof to recognise a line bundle as trivial from its restrictions to the components; the pointer is registered so the reference resolves and the cited statement is not reproduced inside the span. The number is the stable Stacks tag of the target, looked up in the shared tag index.\par\medskip
\ifdefined\hypertarget\hypertarget{sourcepointer-7}{}\fi
\noindent\textit{Source reference.} Original lemma 00GR: focused-seeds/AGP/sources/stacks-snapshot-20261011/tags.txt, source lines 617--617. Allowed input of the Algebra chapter quoted by the retained proof for the assertion that an integral ring extension of a field is a field, which is what identifies the global sections with the residue field; the pointer is registered so the reference resolves and the cited statement is not reproduced inside the span. The number is the stable Stacks tag of the target, looked up in the shared tag index.\par\medskip
\begin{situation}
\label{situation-regular-model}
Let $R$ be a discrete valuation ring with fraction field $K$,
residue field $k$, and uniformizer $\pi$.
Let $C$ be a smooth projective curve over $K$ with $H^0(C, \mathcal{O}_C) = K$.
Let $X$ be a regular proper model of $C$.
Let $C_1, \ldots, C_n$ be the irreducible components of the special
fibre $X_k$. Write $X_k = \sum m_i C_i$ as in
Lemma \ref{lemma-divisor-special-fiber}.
\end{situation}

\begin{lemma}
\label{lemma-regular-model-field}
\begin{reference}
\cite[Lemma 2.6]{Artin-Winters}
\end{reference}
In Situation \ref{situation-regular-model} let $d = \gcd(m_1, \ldots, m_n)$.
Let $D = \sum (m_i/d) C_i$ as an effective Cartier divisor. Then there exists
a sequence of effective Cartier divisors
$$
(X_k)_{red} = Z_0 \subset Z_1 \subset \ldots \subset Z_m = D
$$
such that $Z_j = Z_{j - 1} + C_{i_j}$ for some $i_j \in \{1, \ldots, n\}$
for $j = 1, \ldots, m$ and such that $H^0(Z_j, \mathcal{O}_{Z_j})$
is a field finite over $k$ for $j = 0, \ldots m$.
\end{lemma}

\begin{proof}
The reduction $D_{red} = (X_k)_{red} = \sum C_i$ is connected
(Lemma \ref{lemma-regular-model-connected}) and proper over $k$. Hence
$H^0(D_{red}, \mathcal{O})$ is a field and a finite extension of
$k$ by Varieties, Lemma
\ref{varieties-lemma-proper-geometrically-reduced-global-sections}.
Thus the result for $Z_0 = D_{red} = (X_k)_{red}$ is true.
Suppose that we have already constructed
$$
(X_k)_{red} = Z_0 \subset Z_1 \subset \ldots \subset Z_t \subset D
$$
with $Z_j = Z_{j - 1} + C_{i_j}$ for some $i_j \in \{1, \ldots, n\}$
for $j = 1, \ldots, t$ and such that $H^0(Z_j, \mathcal{O}_{Z_j})$
is a field finite over $k$ for $j = 0, \ldots, t$.
Write $Z_t = \sum a_i C_i$ with $1 \leq a_i \leq m_i/d$.
If $a_i = m_i/d$ for all $i$, then $Z_t = D$ and the lemma is proved.
If not, then $a_i < m_i/d$ for some $i$ and it follows that
$(Z_t \cdot Z_t) < 0$ by Lemma \ref{lemma-properties-form}. This means that
$(D - Z_t \cdot Z_t) > 0$ because $(D \cdot Z_t) = 0$ by the lemma.
Thus we can find an $i$ with $a_i < m_i/d$ such that
$(C_i \cdot Z_t) > 0$. Set $Z_{t + 1} = Z_t + C_i$ and $i_{t + 1} = i$.
Consider the short exact sequence
$$
0 \to \mathcal{O}_X(-Z_t)|_{C_i} \to \mathcal{O}_{Z_{t + 1}} \to
\mathcal{O}_{Z_t} \to 0
$$
of Divisors, Lemma \ref{divisors-lemma-ses-add-divisor}.
By our choice of $i$ we see that
$\mathcal{O}_X(-Z_t)|_{C_i}$ is an invertible sheaf of negative degree
on the proper curve $C_i$, hence it has no nonzero global sections
(Varieties, Lemma \ref{varieties-lemma-check-invertible-sheaf-trivial}).
We conclude that $H^0(\mathcal{O}_{Z_{t + 1}}) \subset H^0(\mathcal{O}_{Z_t})$
is a field (this is clear but also follows from
Algebra, Lemma \ref{algebra-lemma-integral-under-field})
and a finite extension of $k$. Thus we have extended the sequence.
Since the process must stop, for example because $t \leq \sum (m_i/d - 1)$,
this finishes the proof.
\end{proof}

\begin{thebibliography}{10}

\bibitem{Artin-Winters}
M.~Artin and G.~Winters.
\newblock Degenerate fibres and stable reduction of curves.
\newblock {\em Topology}, 10:373--383, 1971.

\end{thebibliography}
\end{document}
